\documentclass[11pt,a4paper]{article}
\usepackage[utf8]{inputenc}
\usepackage[T1]{fontenc}
\usepackage{estilo_capstone}

\titulo{Definition of Done — Evento Cultural}
\subtitulo{Criterios de terminación del trabajo (artefacto 4 del marco ágil, anexo metodológico)}
\autor{Roberto Riffo · José Miguel Moncada · David Herrera}
\fecha{25 de septiembre de 2026}

\newcolumntype{L}[1]{>{\raggedright\arraybackslash}p{#1}}

\begin{document}
\maketitlecapstone

\section{Propósito y alcance}

Este documento define los criterios que debe cumplir una historia de usuario o
tarea del equipo para considerarse \textbf{terminada}. Se elabora a partir de
los criterios de aceptación ya definidos en el backlog de historias de usuario
(\code{Historias\_de\_Usuario\_Plataforma\_Cultural.xlsx}) y en la matriz de
requerimientos adaptada a Fase 2
(\code{Requerimientos\_...\_MoSCoW\_ADAPTADO\_Fase2.xlsx}). No declara
verificaciones que el equipo no haya ejecutado: dónde falta ratificar un
criterio, se marca como pendiente de decisión del equipo.

\section{Criterios de terminación}

\begin{enumerate}
  \item \textbf{Funcionalidad implementada} según la historia o tarea, sin
        funcionalidad inventada: lo declarado como implementado debe existir en
        el código.
  \item \textbf{Datos vigentes:} ninguna vista muestra eventos vencidos; el
        estado del evento se deriva de sus fechas o se declara explícitamente
        (suspendido/cancelado).
  \item \textbf{Origen del dato visible:} todo evento presenta su
        \code{fuente} y \code{ultimaComprobacion}; los enlaces no disponibles
        se muestran como tales, sin anclas rotas.
  \item \textbf{Verificación del frontend aprobada:} \code{ng build
        --configuration production}, \code{npm run lint} y \code{npm test}
        (Vitest) sin errores, y consola del navegador limpia en el recorrido de
        la funcionalidad.
  \item \textbf{Verificación del scraper aprobada:} la prueba unitaria del
        módulo o conector afectado pasa (\code{python -m unittest discover -s
        tests}), y la ejecución completa del pipeline termina sin error,
        dejando registro en la bitácora de ejecuciones.
  \item \textbf{Accesibilidad mínima verificada:} recorrido por teclado,
        foco visible, \code{aria-pressed}/\code{aria-live} donde corresponde y
        áreas táctiles de al menos 44 px, conforme a la verificación previa del
        frontend.
  \item \textbf{Estados de interfaz completos:} carga (skeleton), vacío y
        error con recuperación (``Reintentar'' / ``Limpiar filtros'').
  \item \textbf{Diseño con tokens:} sin colores ni medidas hardcodeadas; el
        estilo consume los tokens primitivos \code{--pr-*} y semánticos
        \code{--ec-*}.
  \item \textbf{Movimiento reducido respetado:} las animaciones nuevas se
        desactivan o reducen con \code{prefers-reduced-motion}.
  \item \textbf{Documentación al día:} si el cambio afecta pantallas,
        servicios o el modelo de datos, la documentación técnica del frontend
        (\code{documentacion-frontend.tex}, control de cambios §14) o la del
        scraper se actualiza en la misma intervención.
  \item \textbf{Historia documentada:} la historia de usuario queda con su
        criterio de aceptación, prioridad MoSCoW y estado de evidencia en el
        backlog; la fila correspondiente del tablero Trello se cierra.
  \item \textbf{Sin secretos ni credenciales} en el repositorio; las
        variables de entorno van documentadas en \code{.env.example}.
\end{enumerate}

\section{Criterios para historias de recomendación (``para ti'')}

Además de los anteriores, una historia del sistema de recomendaciones se
considera terminada cuando: cada recomendación muestra un motivo legible para
la persona (``Porque exploraste\ldots''), el puntaje de afinidad se compone de
las señales registradas en la visita, y ``No me interesa'' retira el evento de
las recomendaciones de forma inmediata.

\section{Pendiente de decisión del equipo}

\begin{pendienteenv}
ratificación de este Definition of Done por los tres integrantes y, si el
equipo lo aprueba, su publicación en el tablero Trello como parte del flujo
Scrumban.
\end{pendienteenv}

\end{document}
